% fixtures.tex — off-board parts for the sa400_led_monitor test run (TEST_PLAN Appendix D):
%   D.1 bus-simulator jig ("fake EI")   D.2 power pigtail for J2   D.3 key plug position
% Build:  lualatex fixtures.tex   (pdf/build_test_plan.py does this and copies the PDF to docs/img/)
\documentclass[border=6pt]{standalone}
\usepackage{fontspec}
\setmainfont{DejaVu Sans}[Path=/usr/share/fonts/truetype/dejavu/,Extension=.ttf,UprightFont=*,BoldFont=*-Bold,ItalicFont=*-Oblique,BoldItalicFont=*-BoldOblique]
\usepackage[american]{circuitikz}
\usetikzlibrary{calc}
\ctikzset{bipoles/length=0.85cm}
\tikzset{every node/.style={font=\footnotesize}, pin/.style={font=\scriptsize, inner sep=1pt, text=black!70},
         note/.style={font=\scriptsize, text=black!70, align=left}, warn/.style={font=\scriptsize, text=red!75!black, align=left}}
\definecolor{ribbon}{gray}{0.80}
\begin{document}
\begin{circuitikz}[x=1cm, y=1cm]
% ===================================================================== D.1 jig
\node[anchor=west, font=\normalsize\bfseries] at (0,2.9) {D.1 Bus-simulator jig (``fake EI'') --- used in Phase 2};
\node[anchor=west, note] at (0,2.35) {Each monitored line is held high by 150~$\Omega$ (as the terminating drive does) and pulled to ground by a switch (as a 7438 in the EI does).\\ The monitor plugs onto the 3463. Ribbon conductor $N$ = J1 pin $N$ when the stripe is on contact 1 --- check with the DMM before trusting the jig.};

% connector + ribbon
\draw[thick, fill=black!75] (0,0.4) rectangle (1.6,-11.3);
\node[font=\scriptsize, text=white, rotate=90] at (0.8,-5.4) {3M 3463-0001 card-edge IDC};
\fill[ribbon] (1.6,0.1) rectangle (3.4,-11.0);
\draw[red, very thick] (1.6,-10.9) -- (3.4,-10.9);
\node[pin, anchor=north west, text=red!80!black] at (1.6,-11.0) {red stripe = conductor 1};
\node[note, anchor=south west] at (1.7,0.45) {$\approx$150 mm ribbon, cut end $\to$};

% signal rows
\foreach \y/\sig/\pin/\n [count=\i] in {-0.6/DS1/10/1, -2.6/DS2/12/2, -4.6/DS3/14/3, -6.6/DS4/6/4, -8.6/WG/24/5} {
  \draw (3.4,\y) -- (5.2,\y) node[circ] (n\n) {};
  \node[pin, anchor=north west] at (3.45,\y-0.04) {\pin\ (\sig)};
  \draw (n\n) to[R, l=150] (7.6,\y) node[circ]{};
  \draw (n\n) -- (5.2,\y-0.6) to[nos, l_=SW \n] (6.7,\y-0.6) -- (6.7,\y-0.95) node[ground, scale=0.8]{};
}
% pull-up rail with jumper
\draw (7.6,-0.6) -- (7.6,-8.6);
\draw (7.6,-0.6) -- (7.6,-0.2) to[nos, l=JP] (7.6,1.1) node[vcc, scale=0.8]{+5V\_JIG};
\node[note, anchor=west] at (8.3,1.5) {bench supply 2 (or the same supply): 5.0 V, 50 mA limit};
\node[note, anchor=west] at (7.95,0.3) {JP closed = pull-ups on (terminated bus)\\ JP open = the board's own 10 k only};
% GND conductors
\draw (3.4,-10.5) -- (5.2,-10.5) -- (5.2,-10.85) node[ground, scale=0.8]{};
\node[pin, anchor=north west] at (3.45,-10.54) {5, 9, 11, 13, 23 (GND) --- twist together};
\node[note, anchor=west] at (5.6,-10.5) {jig GND = bench supply (--) = the monitor supply's (--): one common ground};
% meter note
\node[note, anchor=north west] at (3.4,-11.6) {Steps 2.6.2 and 2.8: mA meter in place of SW\,$n$ (line to GND). Step 2.6.22: meter in series with that line's 150~$\Omega$.\\ Any 34-way 0.1\,in card-edge IDC and 28~AWG ribbon will do; parts in Appendix C.};

\begin{scope}[yshift=-3.3cm]
% ===================================================================== D.2 pigtail
\node[anchor=west, font=\normalsize\bfseries] at (0,-10.2) {D.2 Power pigtail for J2 --- used in Phases 2 and 3};
% supply 1
\draw[thick] (0,-11.0) rectangle (2.6,-13.0);
\node[font=\scriptsize, align=center] at (1.3,-11.75) {bench supply 1\\ \textbf{5.00 V}\\ 100 mA limit};
\node[pin, anchor=west] at (2.62,-11.3) {+};
\node[pin, anchor=west] at (2.62,-12.7) {--};
\draw (2.6,-11.3) -- (3.3,-11.3) to[ammeter, l=mA] (5.0,-11.3) -- (7.6,-11.3);
\draw (2.6,-12.7) -- (7.6,-12.7);
\draw (6.8,-12.7) node[circ]{} -- (6.8,-12.2) -- (7.6,-12.2);
% supply 2
\draw[thick] (0,-13.6) rectangle (2.6,-15.6);
\node[font=\scriptsize, align=center] at (1.3,-14.35) {bench supply 2\\ \textbf{12.0 V}\\ 50 mA limit\\ (optional)};
\node[pin, anchor=west] at (2.62,-13.9) {+};
\node[pin, anchor=west] at (2.62,-15.3) {--};
\draw (2.6,-13.9) -- (7.6,-13.9);
\draw (2.6,-15.3) -- (6.4,-15.3) -- (6.4,-12.7) node[circ]{};
% plug
\draw[thick, fill=black!10] (7.6,-10.9) rectangle (9.8,-14.3);
\node[font=\scriptsize\bfseries, align=center] at (8.7,-11.0) {};
\foreach \y/\p/\nm in {-11.3/4/+5 V, -12.2/3/RTN, -12.7/2/RTN, -13.9/1/+12 V} {
  \draw[fill=white] (7.6,\y) circle (0.12);
  \node[pin, anchor=west] at (7.8,\y) {\p\ \ \nm};
}
\node[note, anchor=north, align=center] at (8.7,-14.4) {AMP 1-480424-0 housing\\ + 60619-1 pins\\ = PC peripheral power plug};
\node[note, anchor=west] at (5.2,-10.95) {red wire};
\node[note, anchor=west] at (3.3,-12.45) {black};
\node[note, anchor=west] at (3.3,-13.65) {yellow};
\node[warn, anchor=north west] at (10.2,-10.9) {Before first use, DMM on the plug:\\ pin 4 = +5 V, pin 1 = +12 V, 2 and 3 = 0 V.\\ Reversed +5 V is blocked by Q1;\\ +12 V on pin 4 is fault F2 and may cost D11.};
\node[note, anchor=north west] at (10.2,-12.9) {Both supplies' (--) are tied at pin 2/3.\\ Keep the mA meter in the +5 V lead for\\ every powered step; it is the board's\\ idle-current fingerprint (2.1.1).};

% ===================================================================== D.3 key plug
\node[anchor=west, font=\normalsize\bfseries] at (0,-17.0) {D.3 Key plug (optional) --- 3M 3439-2 in the connector, opposite the board's slot between fingers 4 and 6};
% board tab, component side, fingers pointing up, pin 2 at right
\draw[thick] (0.5,-19.6) -- (0.5,-18.0) -- (9.5,-18.0) -- (9.5,-19.6);
\foreach \i in {0,...,16} {
  \pgfmathsetmacro{\x}{9.1 - \i*0.5}
  \pgfmathtruncatemacro{\p}{2 + 2*\i}
  \pgfmathsetmacro{\w}{ifthenelse(\i==1 || \i==2, 0.08, 0.12)}
  \draw[fill=yellow!60] (\x-\w,-19.4) rectangle (\x+\w,-18.2);
  \node[pin, anchor=north] at (\x,-19.55) {\p};
}
\draw[fill=white, thick] (8.3,-18.0) rectangle (8.4,-19.6);
\node[note, anchor=south] at (8.35,-17.95) {slot};
\node[note, anchor=north west] at (0.5,-19.9) {Board tab seen from the component side, fingers pointing away: pin 2 at the right, even pins on top, odd (GND) underneath. The 3439-2 pushes into the 3463's\\ contact cavity at that position; the connector then only fits the monitor one way. Without it, the stripe-to-pin-1 check in 3.0 is what prevents a reversed plug.};
\end{scope}
\end{circuitikz}
\end{document}
